\documentclass[10pt,a4paper]{article}
\usepackage[margin=2.5cm]{geometry}
\usepackage[utf8]{inputenc}\usepackage[T1]{fontenc}
\usepackage[spanish,es-nodecimaldot]{babel}
\usepackage{amsmath,amssymb,amsthm,booktabs}
\usepackage[colorlinks=true,urlcolor=blue]{hyperref}
\newtheorem{proposition}{Proposición}
\title{Escalado de acoplamiento con bloques de incidencia\\Extensión espectral acotada}
\author{Nota separada del manuscrito OMS}\date{18 de octubre de 2026}
\begin{document}\maketitle
\noindent\textbf{Alcance.} Se obtiene una fórmula cerrada para el mejor factor común
en una subclase con capacidades uniformes, sin blanquear los bloques locales ni
suponer rango completo. Para capacidades heterogéneas se obtiene una condición
suficiente verificable. Son resultados de condicionamiento, no mejoras de tiempo,
precisión de puntos u optimalidad MILP. No se ha establecido prioridad bibliográfica.

\section{Descomposición exacta y ceros}
Sean $B\in\mathbb R^{r\times n}$, $C\in\mathbb R^{s\times n}$,
$D\in\mathbb R^{s\times p}$, $K\ge1$, y $\Gamma$ diagonal positiva. Considerar
\[
 A_\Gamma=\begin{bmatrix}I_K\otimes B&0\\
 \mathbf1_K^{\mathsf T}\otimes\Gamma C&\Gamma D\end{bmatrix},
 \qquad
 M_\Gamma=\begin{bmatrix}B&0\\\sqrt K\Gamma C&\Gamma D\end{bmatrix}.
\]
\begin{proposition}
Los valores singulares positivos de $A_\Gamma$, con multiplicidades, son los de
$M_\Gamma$ y $K-1$ copias de los de $B$. Incluidos todos los ceros, la identidad
inequívoca es
\[
 \operatorname{spec}(A_\Gamma^{\mathsf T}A_\Gamma)=
 \operatorname{spec}(M_\Gamma^{\mathsf T}M_\Gamma)
 \uplus (K-1)\operatorname{spec}(B^{\mathsf T}B).
\]
\end{proposition}
\begin{proof}
Elegir una matriz ortogonal $Q$ cuya primera columna sea
$\mathbf1_K/\sqrt K$. Multiplicar las filas locales por
$Q^{\mathsf T}\otimes I_r$ y las columnas locales por $Q\otimes I_n$.
El bloque local sigue siendo $I_K\otimes B$, mientras el acoplamiento se convierte
en $[\sqrt K\Gamma C,0,\ldots,0]$. Permutar filas y columnas da la suma directa
rectangular de $M_\Gamma$ y las $K-1$ copias de $B$.
\end{proof}
La unión de las SVD económicas puede omitir ceros cuando los bloques tienen
distintas orientaciones rectangulares. Por ejemplo, $K=2$, $B=[1\ 0]$,
$C=[I_2;0\ 0]$ y $D=[0\ 0\ 1]^{\mathsf T}$ dan una matriz $A$ de $5\times5$
con rango cuatro: su SVD tiene cinco valores, pero las SVD económicas de $M$ y
$B$ aportan solo cuatro. La identidad sobre valores positivos permanece válida.

\section{Mejor factor común con capacidad uniforme}
Ahora $C=I_n$, $D=dI_n$, $d>0$, $\Gamma=\gamma I_n$, $\gamma>0$, y
$B\ne0$ puede ser rectangular y deficiente de rango. Sean
$a=\sigma_{\min}^+(B)$ y $b=\sigma_{\max}(B)$.
\begin{proposition}
El único factor común que minimiza $\kappa_2^+(A_\gamma)$ es
\[
 \boxed{\gamma_* =\sqrt{\frac{ab}{K+d^2}}.}
\]
Con $\theta=d^2/(K+d^2)$ y $q=b/a$, el mínimo es
\[
 \kappa_* = \frac{1+q+\sqrt{(1+q)^2-4\theta q}}{2\sqrt\theta}.
\]
Sobre $0<L\le\gamma\le H$, el óptimo es $\min(H,\max(L,\gamma_*))$.
\end{proposition}
\begin{proof}
Una SVD de $B$ reduce el problema a bloques de dos dimensiones. Escribir
$t=(K+d^2)\gamma^2$. Para cada valor singular $s>0$, los cuadrados de los
valores singulares correspondientes son
\[
 \lambda_\pm(s^2,t)=\frac{s^2+t\pm\sqrt{(s^2+t)^2-4\theta s^2t}}2.
\]
Los modos del núcleo de $B$ aportan valores positivos $\sqrt t$ y ceros.
La mayor raíz aumenta con $s$, la menor también aumenta y es menor que $s^2$;
además $\lambda_+(s^2,t)\ge t$ y $\lambda_-(s^2,t)\le\theta t$.
Por tanto, incluidos los contrastes de mercancías,
\[
 \kappa_2^+(A_\gamma)^2
 =\frac{\lambda_+(b^2,t)}{\lambda_-(a^2,t)}
 =\frac{\lambda_+(b^2,t)\lambda_+(a^2,t)}{\theta a^2t}.
\]
Definir $h(u)=\tfrac12\log\theta+
\operatorname{arcosh}(\cosh(u/2)/\sqrt\theta)$. Entonces
$\log\lambda_+(s^2,t)=\tfrac12\log(s^2t)+h(\log t-\log s^2)$.
La función $h$ es par y estrictamente convexa, pues $0<\theta<1$ y
\[
 h''(u)=\frac{(1-\theta)\cosh(u/2)}{4(\cosh^2(u/2)-\theta)^{3/2}}>0.
\]
El objetivo logarítmico es una constante más
$h(\log t-\log a^2)+h(\log t-\log b^2)$. Su mínimo único está en
$\log t=\tfrac12(\log a^2+\log b^2)$, es decir, $t=ab$.
Sustituir da ambas fórmulas. La convexidad demuestra la proyección sobre el intervalo.
\end{proof}
El teorema admite directamente matrices de incidencia. La capacidad uniforme
es una hipótesis real de esta fórmula; no se atribuye su optimalidad a matrices
$D$ heterogéneas ni a factores diagonales independientes.

\section{Capacidades heterogéneas: condición suficiente}
Sea ahora $C=I_n$, $\Gamma=\gamma I_n$ y
$0<d_-^2I\preceq DD^{\mathsf T}\preceq d_+^2I$.
Definir
\[
 \Lambda_+(s^2,d,\gamma)=\frac{s^2+\gamma^2(K+d^2)+
 \sqrt{(s^2+\gamma^2(K+d^2))^2-4s^2d^2\gamma^2}}2.
\]
El orden de Loewner entre los Gram de filas, junto con la sección anterior,
proporciona
\[
 \kappa_2^+(A_\gamma)^2\le
 U(\gamma):=\frac{\Lambda_+(b^2,d_+,\gamma)
 \Lambda_+(a^2,d_-,\gamma)}{a^2d_-^2\gamma^2}.
\]
Los núcleos de esos Gram coinciden: $D$ tiene rango completo de filas, por lo
que cualquier vector del núcleo tiene componente de acoplamiento cero y
componente local en $\ker B^{\mathsf T}$. Por ello el orden se aplica también
a los menores autovalores positivos. Si se verifica
$U(\gamma)<\kappa_2^+(A_1)^2$, se certifica una mejora estricta.
Una cota inferior rigurosa para el segundo miembro basta para comprobarlo.

\section{Certificados sobre las seis R nativas}
En la formulación dirigida R, $D=-\operatorname{diag}(U_e)$ y el grafo es
conexo. Se aislaron exactamente las raíces del polinomio entero de $BB^{\mathsf T}$
para acotar $a^2$ y $b^2$. Se usaron cotas racionales superiores de las raíces
cuadradas en $U(\gamma)$; no se infiere el certificado de una SVD flotante.
Para el modelo sin escalar,
\[
 \kappa_2^+(A_1)^2\ge\frac{K+U_{\max}^2}{H},\qquad
 H=\min_{u<v}\frac{L_{uu}+L_{vv}-2L_{uv}}2,\quad L=BB^{\mathsf T}.
\]
La fila de máxima capacidad acota el máximo autovalor por debajo; un vector
local $e_u-e_v$, ortogonal al núcleo, acota el mínimo positivo por arriba.

Los factores son diádicos y cumplen en todas las capacidades
$10^{-6}/\gamma+U_e10^{-9}\le10^{-5}$, además de la ventana de coeficientes.
Se verifica la desigualdad estricta entre cotas racionales en los seis casos.
Las cifras siguientes son aproximaciones para lectura; las fracciones exactas
se conservan en \texttt{results.json}.
\begin{center}\small
\begin{tabular}{lrrr}\toprule
Caso&$\gamma$&Cota inferior de $\kappa(A_1)$&Cota superior de $\kappa(A_\gamma)$\\\midrule
R01.1&$1/8$&274.15&40.09\\
R02.1&$1/8$&633.71&103.74\\
R03.1&$1/4$&1125.70&335.16\\
R04.1&$1/8$&198.89&26.91\\
R05.1&$1/8$&459.74&60.70\\
R06.1&$1/4$&856.53&233.50\\\bottomrule
\end{tabular}
\end{center}
Esto certifica el condicionamiento de la matriz completa en aritmética real.
No es una cota de error de una solución ni del trabajo del optimizador. El
experimento previo sobre estas mismas redes ya mostró que mejores matrices no
implican una ventaja práctica. Los certificados no sustituyen esa evidencia.

\section{Comprobaciones y decisión}
Se verifican 40 identidades de espectros Gram, un contraejemplo para el conteo de
ceros, 135 comparaciones de la fórmula con SVD densas y seis certificados nativos
con racionales e intervalos algebraicos exactos. No se ejecutó otro MIP para esta nota.
Se conserva separada del artículo. La optimización diagonal y por bloques tiene
antecedentes clásicos y modernos; esta búsqueda acotada no acredita novedad de
la fórmula.\footnote{\url{https://web.stanford.edu/~boyd/lmibook/lmibook.pdf}, capítulo 3;
\url{https://arxiv.org/abs/2209.00809}.}
\end{document}
